
%----------------------------------------
\begin{frame}[fragile]
  \frametitle{Conclusion} \small

\begin{itemize}

  \item Temporal logic can be a powerful formalism for expressing
  temporal constraints on sequences, here used to {\bf generate} sequences.

  \pause

  \item One can develop alternative (perhaps more userfriendly) formalisms
  for expressing temporal constraints, including e.g graphical formalisms.

  \pause

  \item Constraint solving using SMT solvers (such as Z3) is an efficient
  and automated way of generating sequences (tests, schedules, ...) satisfying temporal constraints.

  \pause

  \item Disclaimer: scalability can, as usual for these kinds of techniques, be an issue. 

  \pause

  \item Temporal properties can also be used to {\bf monitor} sequences. This is
  the field of {\em runtime verification}. That's another talk.

\end{itemize}

%\begin{block}{fuzz is available here}
%\begin{center}
%\url{https://github.com/havelund/fuzzing}
%\end{center}
%\end{block}

\end{frame}
